\documentclass[11pt,a4paper]{article}
\input{gift_preamble}
% Essay layout: unnumbered sections (headings carry their own numbers).
\setcounter{secnumdepth}{0}
\emergencystretch=3em
\fancyhead[L]{The Constants of Nature Are Counts}
\hypersetup{pdftitle={The Constants of Nature Are Counts: a testable geometric hypothesis for the Standard Model}}

\title{%
\LARGE\textbf{The Constants of Nature Are Counts}\\[0.45em]
\large A testable geometric hypothesis for the Standard Model.\\[0.15em]
\normalsize What is derived. What is assumed. What would refute it.%
}
\author{}
\date{}

\begin{document}

\maketitle
\noindent\rule{\textwidth}{0.2pt}

\noindent\textbf{Author}: Brieuc de La Fourni\`ere

\noindent Independent researcher, Beaune, France \quad | \quad Arithmon program

\noindent Contact: \href{mailto:brieuc@arithmon.com}{\texttt{brieuc@arithmon.com}} \quad | \quad \href{https://arithmon.com}{arithmon.com}

\noindent\textbf{Date}: July 2026

\noindent\rule{\textwidth}{0.2pt}

\begin{quote}
\textbf{Written for skeptical physicists.} If after fifteen minutes you
conclude this is wrong, it should be clear \emph{where} it is wrong. That
is the standard this document is trying to meet.
\end{quote}

\section{Why read this}

Every successful physical theory eventually explains numbers that once had
to be measured. The Standard Model has not yet done so. It predicts
thousands of measurements to extraordinary precision. And it leaves one
question almost entirely unanswered: why do the fundamental constants have
the values they do?

Why these particle masses. Why these mixing angles. Why exactly three
generations. Today those nineteen numbers (twenty-six, once neutrino masses
and mixings are counted) are measured and inserted by hand.

The $\Kseven$ framework explores a single possibility: that none of them
are free, and that they are instead arithmetic consequences of one compact
seven-dimensional geometry. A hypothesis of that shape has only two
available fates. It is spectacularly wrong, or it leaves measurable traces
everywhere. There is no comfortable middle in which it is vaguely
suggestive.

What follows is the claim, the assumptions it rests on, the evidence, and
the specific ways it can fail. If correct, the framework changes the status
of the Standard Model parameters from empirical inputs to geometric
outputs.

\section{At a glance}

\begin{center}
\begin{tabular}{@{}p{0.34\textwidth}p{0.60\textwidth}@{}}
\toprule
Core hypothesis & Standard Model constants are invariants of one compact $\Gtwo$ geometry \\
Continuously adjustable parameters & \textbf{0}, after the input ledger is frozen \\
Parameter-free algebraic relations (Type~\typeI) & \textbf{33} \\
Total derived observables, all types & \textbf{95} (66 experimentally comparable) \\
Agreement & \textbf{11} exact to better than 0.01\%, \textbf{53} within 1\% \\
Formal verification & \textbf{213}-conjunct certificate in Lean~4, \textbf{15} classified axioms, \textbf{0} unproven steps \\
Pre-registration & ledger, grammar and target list deposited before any search (\href{https://doi.org/10.5281/zenodo.20666879}{DOI}) \\
\textbf{Falsification} & \textbf{$\delta_{\CP}$ measured outside $[182°, 212°]$ at $3\sigma$ and the framework is wrong. DUNE beam targeted 2031.} \\
\bottomrule
\end{tabular}
\end{center}

That last line is the one that matters most, and section~8 returns to it.
A framework with no continuously adjustable parameter cannot absorb a bad
measurement. That is a cost, not a feature to be apologised for, and it is
what makes the rest of the table worth checking.

\section{1. The claim}

The dimensionless constants of the Standard Model are arithmetic and
topological invariants of a compact seven-dimensional manifold with
$\Gtwo$ holonomy, not free parameters to be fitted.

That is the whole hypothesis. The specific realisation is a geometry with
Betti numbers $(\btwo, \bthree) = (21, 77)$ and an $\E_8 \times \E_8$
motivated exceptional architecture. From it, thirty-three exact algebraic
relations among Standard Model and cosmological observables follow with no
continuously adjustable quantity remaining once the geometric input is
declared.

\section{2. Why this deserves skepticism}

If someone claims to derive Nature's constants, your first reaction should
be skepticism. Mine would be too. Any sufficiently large pool of small
integers and transcendentals produces impressive coincidences, and the
history of this genre is not encouraging.

There are three obvious failure modes. The framework's response to each is
a commitment made in advance rather than an argument made afterwards.

\begin{center}
\begin{tabular}{@{}p{0.30\textwidth}p{0.64\textwidth}@{}}
\toprule
\textbf{Objection} & \textbf{Commitment made before any result was obtained} \\
\midrule
This is numerology & Ledger, formula grammar and target list frozen and publicly deposited before the search; four independent null models \\
The relations were cherry-picked & The formula grammar was declared in advance, so the search space is auditable \\
There are hidden fitted parameters & The ledger is explicit and finite; no continuous adjustment is available after freeze \\
\bottomrule
\end{tabular}
\end{center}

None of this proves the framework correct. It makes the standard
objections testable instead of rhetorical, which is a lower bar and the
only one that can be cleared in advance.

The methodology behind the null models is deliberately not bespoke. It is
calibrated against historical verdicts, so that Eddington's fine-structure
argument fails it and the quantum Hall relation passes it. A test that only
its author's framework can pass is not a test.

\section{3. What is assumed}

The inputs are a declared ledger of twenty structural constants, all
algebraic functions of six primitive topological integers: $\btwo$,
$\bthree$, $\dim(\Gtwo)$, $\dim(\E_8)$, $\rank(\E_8)$ and
$\dim(\Kseven)$. Standard transcendentals enter as well ($\pi$, $\sqrt{2}$,
$\ln 2$, $\zeta$ values, and the golden ratio $\varphi$, whose appearance
is traced to the McKay correspondence $\E_8 \leftrightarrow 2I$ and carries
its own caveat).

Some repository summaries compress the same input further, to three deeper
integer choices: $N_{\rm gen} = 3$, $\rank(\E_8) = 8$, $\rank(\Gtwo) = 2$.
The two counts describe one structure. Naming the groups fixes
$\dim(\E_8)$, $\dim(\Gtwo)$ and $\dim(\Kseven)$, so the three choices
determine four of the six integers. What they do not determine is
$(\btwo, \bthree)$: that selection is the open question of section~8,
and no compression hides it.

The ledger, the grammar and the target list were frozen and deposited
\textbf{before} any relation was searched for. This matters more than any
individual result below. Pre-registration is what separates a prediction
from a retrofit, and here it rests on a timestamp rather than on the
author's assurance.

One assumption should be handed to a critical reader rather than left to
be discovered. The metric normalisation $\det(g) = 65/32$ is imposed as a
target, not derived. Six observables depend on it. Its expression in terms
of topological integers is suggestive but not derivational. If there is a
soft joint in the construction, it is this one.

\section{4. What is derived}

A recurring problem with ambitious frameworks is that everything
eventually gets called a prediction. $\Kseven$ separates four epistemic
categories and keeps them separate in every table it publishes.

\begin{center}
\begin{tabular}{@{}clp{0.52\textwidth}l@{}}
\toprule
\textbf{Type} & \textbf{Count} & \textbf{Nature} & \textbf{Mean deviation} \\
\midrule
\typeI & 33 & Direct algebraic identities in the frozen ledger & 0.73\% (frozen dataset) \\
\typeII & 19 & One-step extractions using a named experimental anchor & 0.17\% \\
\typeIII & 21 & Multi-step dynamical chains, conditional on the mechanism invoked & 3.4\% \\
\typeIV & 22 & Structural diagnostics, no direct experimental comparison & n/a \\
\bottomrule
\end{tabular}
\end{center}

One bookkeeping point, disclosed rather than discovered: two Type~\typeI{}
deviation figures circulate and both are real. The 0.73\% above is the
mean against the dataset frozen at pre-registration (PDG 2024 + NuFIT
6.0), which is the figure this document quotes because it is the one the
pre-registration protects. After a post-freeze audit reconciled six
experimental values to primary sources, the same mean reads 0.99\%, and
that is the figure the repository summaries carry. Supplement S4 documents
both.

Only the thirty-three Type~\typeI{} relations are parameter-free
consequences of the frozen ledger. The nineteen Type~\typeII{} relations
are conditional reconstructions: each multiplies a Type~\typeI{} ratio by
a measured anchor, so that $m_u$ is obtained as $(m_u/m_d) \times m_d$
measured. They are nineteen traceable extractions of physical scale, not
nineteen independent predictions, and counting them as the latter would be
the easiest available way to oversell this work. The twenty-one
Type~\typeIII{} chains are conditional on the mechanism each invokes,
enumerated case by case. Type~\typeIV{} are diagnostics.

Of the sixty-six experimentally comparable observables, eleven agree to
better than 0.01\% and fifty-three fall within 1\%.

\section{5. Why the numbers are not the main result}

Numerical precision is reported here as a secondary figure. That is a
deliberate inversion of the usual presentation, and it is not modesty.

A framework of this kind can always be made to look impressive by leading
with sub-percent agreement across dozens of quantities. But precision is
cheap to manufacture and expensive to interpret: it is exactly what a
well-tuned fit also produces. The claim being made is structural, namely
that these relations are algebraic consequences of a geometric input fixed
in advance. Precision is evidence about that claim. It is not the claim.

An earlier version of this work led with joint coincidence probabilities
of order $10^{-346}$ against a uniform null and $10^{-133}$ against an
algebraic null over 4.2 million random formulas. Those figures are real,
they are retained in the supplements as internal-consistency diagnostics,
and they have been deliberately removed from the headline. They measure
something, but not what a skeptical reader needs to know, and presenting
them as primary evidence was a methodological error this version corrects.

\section{6. Formal verification, and its limits}

The framework carries a Lean~4 formal layer. At current head, the master
certificate spans ten files with a top-level conjunct count of
\textbf{213}, resting on \textbf{15 classified axioms} in an A to F
taxonomy, with \textbf{zero} unproven steps. Four of the fifteen axioms
are external data packages, named in the source:
\texttt{K7\_analysis\_data}, \texttt{K7\_spectral\_data},
\texttt{literature\_package}, \texttt{KK\_YM\_EFT}. Fifty-five of the
ninety-five observables fall inside the verified perimeter.

What this establishes: the algebraic relations hold as stated, and the
internal consistency of the ledger is not a matter of arithmetic trust.

What it does not establish, stated plainly because formal verification is
routinely oversold: Lean certifies no physical interpretation, and it does
not certify the existence of the complete compact construction. A
machine-checked identity between topological integers is exactly that. The
certificate drives the probability of a bookkeeping error to zero and
leaves every substantive question open.

\section{7. Why this is not numerology, as a testable proposition}

Applied to the thirty-three Type~\typeI{} relations, the pre-registered
procedure runs a four-null battery (uniform, algebraic, factorised,
permuted) through the same declared twenty-constant ledger, then ranks
survivors by budget uniqueness.

The outcome is modest and is reported as such. The unique survivor at
exact rank-one budget uniqueness is $m_H/m_W$. Koide's $Q = 2/3$ survives
with a narrow margin. That is the honest statistical claim: two relations
distinguish themselves under a null battery fixed in advance. The
thirty-three as a set are conditional algebraic identities of the frozen
ledger, machine-checked, but not individually distinguished by the sieve.

Two further checks target the most likely failure mode, which is
unconscious selection rather than deliberate fraud.

\textbf{Does formula complexity buy precision?} A random-forest regression
predicting absolute deviation from formula features achieves
$R^2 = -0.518$ under leave-one-out cross-validation, which is worse than a
mean predictor. Complex formulas do not outperform simple ones.

\textbf{Are the good agreements simply the insensitive ones?} Perturbing
each of the twenty structural constants by $\pm 1$ gives a $20 \times 33$
sensitivity matrix. The correlation between sensitivity and deviation is
$\rho = -0.083$.

Neither result proves the framework correct. Both would look very
different had the relations been selected by fitting.

\section{8. The three questions that matter most}

These are load-bearing and unresolved. They appear here at full
prominence, because a reader who finds them on page sixty of the full
paper is entitled to feel handled.

\textbf{Does the geometry exist?} The framework needs a global compact
torsion-free $\Gtwo$ structure on a smooth $\Kseven$. The datum-level
analytic layer is discharged conditionally in a companion paper. The
global-analytic layer remains conditional on a two-slot hypothesis pack.
The framework does not claim the construction exists; it claims what
follows if it does, and it labels which layer each result sits above.

\textbf{Where does the gauge sector come from?} The branching chain
$\E_8 \supset \E_6 \times \SU(3) \supset \cdots \supset \SU(3) \times
\SU(2) \times \U(1)$ is compatible with Standard Model matter content, but
it is a chain of subgroup branchings, not a physical breaking mechanism.
The bridge to a non-abelian chiral gauge sector is open. This is the
largest gap in the work.

\textbf{Why $(21, 77)$?} The selection question was audited, not answered.
Five candidate routes to deriving the Betti numbers were investigated and
closed. One explicit residual remains: $\btwo + \bthree = 98 =
\dim(\Kseven) \cdot \dim(\Gtwo)$. No uniqueness is claimed, either of the
geometry or of the assignment of formulas to observables.

\section{9. How this framework could fail}

\begin{center}
\small
\begin{tabular}{@{}lp{0.27\textwidth}p{0.145\textwidth}p{0.135\textwidth}p{0.21\textwidth}@{}}
\toprule
\textbf{Prediction} & \textbf{Value} & \textbf{Experiment} & \textbf{Timeline} & \textbf{Refuted if} \\
\midrule
$\delta_{\CP}$ & $197° = \dim(\Kseven)\cdot\dim(\Gtwo) + H^* = 7 \times 14 + 99$ & DUNE Phase I ($3\sigma$) & beam targeted 2031 & measured outside $[182°, 212°]$ \\
$\delta_{\CP}$ & as above & DUNE Phase II ($5\sigma$) & late 2030s--2040s & definitive \\
$N_{\rm gen}$ & 3 & any & ongoing & a fourth generation \\
$m_s/m_d$ & 20 & Lattice QCD & 2028--2030 & target $\pm 0.5$; current $20.0 \pm 1.0$ \\
$\sin^2\theta_W$ & $3/13$ & FCC-ee & 2040s & precision $\sim 10^{-5}$ \\
\bottomrule
\end{tabular}
\end{center}

Every integer in the $\delta_{\CP}$ decomposition is a named invariant,
because in this framework it has to be: $7 \times 14 = \dim(\Kseven) \cdot
\dim(\Gtwo) = 98$, the same product that reappears as $\btwo + \bthree =
98$ in the residual of section~8, and $99 = H^*$, the effective
cohomological dimension $1 + \btwo + \bthree$.

No reinterpretation is available. If DUNE returns $\delta_{\CP} = 250°$,
this framework is wrong, nothing in it can be tuned to survive, and the
correct response will be to say so in public.

Two further consequences are recorded as complementary rather than
critical, being model-dependent: a proton lifetime of $4.06 \times
10^{38}$ years, beyond near-term Hyper-Kamiokande reach, and a SUSY
spectrum with $m_{\rm gravitino} = 166$ GeV and $m_{\rm moduli} = 3.2$
TeV, viable only in compressed or suppressed-coupling realisations.

\section{10. Check it yourself}

The framework's central methodological commitment is that its claims
should be verifiable without trusting the author, so here is the shortest
path to testing one of them.

The Lean sources are public. The conjunct count in section~6 is the number
of top-level $\wedge$ operands across the master certificate files. Clone
the repository, strip comments, count operands at parenthesis depth zero,
and compare. If the number is not 213, that discrepancy is a fact about
this work and should be reported as one.

The same holds for pre-registration. The deposit predates the search, the
timestamp is on the DOI, and the grammar is in the repository. None of it
requires taking anyone's word.

\section{11. Where to go from here}

\begin{itemize}\tightlist
\item \textbf{The full framework paper}, with four supplements:
  derivations, complete observable tables, axiom accounting, selection
  audit.
\item \textbf{The methodology paper}: the four-null battery, the
  historical calibration, and a scorecard reusable on any framework,
  including this one.
\item \textbf{The formal core}: Lean~4 sources, certificates, build
  instructions.
\item \textbf{The confrontations scoreboard}: named predictions against
  scheduled data, dated, with outcomes recorded as they arrive.
\item \textbf{The program charter}: hard core, heuristics, fair-play
  rules, open problems.
\end{itemize}

At \href{https://arithmon.com}{arithmon.com} and
\href{https://github.com/Arithmon}{github.com/Arithmon}.

\section{The question}

Most proposals that attempt to explain the constants of Nature eventually
fail. Perhaps this one will too.

But the question does not go away with them. If the constants are not
fundamental, they come from somewhere. $\Kseven$ is one attempt to answer
that while making every assumption explicit, every calculation
reproducible, every prediction public, and every route to refutation
available in advance.

Whether the hypothesis survives is now a matter for mathematics and
experiment, not for authority.

\vspace{1.5em}
\noindent\rule{\textwidth}{0.2pt}

{\footnotesize\noindent The underlying numerical work on $\Gtwo$ metrics
is cited as reference 25 in \emph{Physics Letters B} \textbf{878} (2026)
140566 (Heyes, Hirst, S\'a Earp and Silva). The author participated
remotely in the DANGER workshop at the Banff International Research
Station, April 2026, and has no formal training in physics or mathematics:
every claim above is built to be checkable without reference to its
author.\par}

\end{document}
